\subsection{Divide-and-Conquer Algorithm Design}

The algorithm follows the divide-and-conquer
approach described by Bentley \cite{bentley1984}
and Cormen et al.\ \cite{clrs2009}.

Given the subsequence $x_{lo}, \ldots, x_{hi}$,
let $mid = \lfloor (lo + hi)/2 \rfloor$.
Every period $[l, r]$ inside $[lo, hi]$ falls
into exactly one of three cases:

\begin{enumerate}
    \item it lies entirely in the left half,
    $lo \le l \le r \le mid$;

    \item it lies entirely in the right half,
    $mid + 1 \le l \le r \le hi$;

    \item it crosses the midpoint,
    $l \le mid < r$.
\end{enumerate}

The first two cases are solved recursively.
The third case is solved directly in linear
time by the subroutine \textsc{Max-Crossing},
and the best of the three answers is returned.

\subsubsection{Algorithm Pseudocode}

\begin{verbatim}
MAX-CROSSING(X, lo, mid, hi)

1.  leftSum = -infinity;  sum = 0
2.  for i = mid downto lo:
3.      sum = sum + X[i]
4.      if sum > leftSum:
5.          leftSum = sum;  bestL = i
6.  rightSum = -infinity;  sum = 0
7.  for j = mid+1 to hi:
8.      sum = sum + X[j]
9.      if sum > rightSum:
10.         rightSum = sum;  bestR = j
11. return (bestL, bestR,
            leftSum + rightSum)

MAX-SUBARRAY(X, lo, hi)

1.  if lo == hi:
2.      return (lo, hi, X[lo])
3.  mid = floor((lo + hi) / 2)
4.  L = MAX-SUBARRAY(X, lo, mid)
5.  R = MAX-SUBARRAY(X, mid+1, hi)
6.  C = MAX-CROSSING(X, lo, mid, hi)
7.  return the one of L, R, C with
       the largest sum
       (ties: prefer L, then R)
\end{verbatim}

The initial call is
\textsc{Max-Subarray}$(X, 1, n)$.

\subsubsection{Divide, Conquer, and Combine}

\begin{itemize}
    \item \textbf{Divide:} split the horizon at
    its middle minute, in $O(1)$ time.

    \item \textbf{Conquer:} recursively find the
    worst episode within each half.

    \item \textbf{Combine:} find the worst episode
    that spans the split point and return the
    best of the three candidates, in $O(n)$ time.
\end{itemize}

\subsubsection{Explanation in the Traffic Domain}

The traffic engineer begins with the full record
of per-minute net queue changes for an approach.

Rather than examining every possible start and
end minute, which would require considering
$n(n+1)/2$ periods, the record is split into an
earlier half and a later half.

The worst congestion episode is either completely
in the earlier half, completely in the later half,
or it straddles the boundary between them.
The first two possibilities are the same question
asked about a shorter record, so the same procedure
is applied to each half until single minutes remain.

For an episode that straddles the boundary, the
engineer walks backward from the boundary and
records the start minute that accumulates the most
queue growth up to the boundary, then walks forward
from the boundary and records the end minute that
accumulates the most growth after it. Joining these
two pieces gives the worst straddling episode.

For example, consider the following ten minutes of
detector counts on one approach:

\begin{center}
\begin{tabular}{|c|c|c|c|}
\hline
Minute & Arrivals & Departures & Net $x_i$ \\
\hline
1 & 12 & 9 & $+3$ \\
2 & 8 & 10 & $-2$ \\
3 & 14 & 9 & $+5$ \\
4 & 10 & 11 & $-1$ \\
5 & 6 & 12 & $-6$ \\
6 & 13 & 9 & $+4$ \\
7 & 16 & 9 & $+7$ \\
8 & 7 & 10 & $-3$ \\
9 & 15 & 9 & $+6$ \\
10 & 4 & 12 & $-8$ \\
\hline
\end{tabular}
\end{center}

The record is split after minute 5.

\begin{itemize}
    \item The worst episode in minutes 1--5 is
    minutes 1--3, with growth $3 - 2 + 5 = 6$.

    \item The worst episode in minutes 6--10 is
    minutes 6--9, with growth
    $4 + 7 - 3 + 6 = 14$.

    \item For a straddling episode, the best start
    walking backward from minute 5 is minute 1
    (suffix sums $-6, -7, -2, -4, -1$), and the best
    end walking forward from minute 6 is minute 9
    (prefix sums $4, 11, 8, 14, 6$). The straddling
    episode, minutes 1--9, has growth $-1 + 14 = 13$.
\end{itemize}

The algorithm returns minutes 6--9: the queue on
this approach grew by 14 vehicles over those four
minutes, even though it partially drained during
minute 8. This is the period an engineer would
examine first when retiming the signal.

Note that the maximum subarray problem can also be
solved in $O(n)$ time by a scan-based method
\cite{bentley1984}. The divide-and-conquer
algorithm is studied here because it is the
technique required for this section; it also has
the practical advantage that the two halves are
independent and can be processed in parallel,
for example when detector records for different
days are stored on different machines.
